\begin{figure}[H]
\centering
\begin{tikzpicture}[>=Latex,x=1.02cm,y=0.82cm]
    \def\ang{20}
    \clip (-5.05,-0.55) rectangle (5.25,6.25);
    \fill[minkfieldred]
        (\ang:-5.8) -- (\ang:5.8) -- ++(90-\ang:7.0)
        -- ($(\ang:-5.8)+(90-\ang:7.0)$) -- cycle;
    \foreach \x in {-4,-2,0,2,4}
        \draw[mink grid x] (\x,-0.35) -- (\x,6.1);
    \foreach \y in {0,1,...,6}
        \draw[mink grid t] (-5.0,\y) -- (5.1,\y);
    \foreach \i in {-6,-5,...,6}
        \draw[mink grid prime] (\ang:{\i*0.8}) -- ++(90-\ang:7.0);

    \draw[mink axis] (-4.95,0) -- (5.1,0) node[below] {$x$};
    \draw[mink axis] (0,-0.4) -- (0,6.05) node[left] {$ct$};
    \draw[mink axis prime] (\ang:-4.9) -- (\ang:5.0) node[below right] {$x'$};
    \draw[mink axis prime] (90-\ang:-0.5) -- (90-\ang:6.15) node[right] {$ct'$};

    \fill[minklightred] (-4.75,1.05) -- (4.75,4.51) -- (4.75,5.05)
        -- (-4.75,1.59) -- cycle;
    \draw[minkdarkred,very thick] (-4.75,1.32) -- (4.75,4.78);
    \node[minkdarkred,font=\scriptsize,above left=1pt] at (4.35,4.63)
        {$t'=\mathrm{cte}$};

    \foreach \s in {-2.8,0,2.8}{
        \draw[minkdarkred!55,thick] (\s-1.7,-0.4) -- (\s+0.65,6.05);
    }
    \foreach \x/\y in {-3/1.96,0/3.05,3/4.14}{
        \node[draw=minkdarkred,fill=white,line width=0.8pt,
            rounded corners=1.5pt,inner xsep=3pt,inner ysep=1.5pt,
            font=\scriptsize,text=minkdarkred] at (\x,\y)
            {$12{:}00$};
    }

    \draw[black,mink dashed,very thick] (-4.85,3.05) -- (4.95,3.05);
    \node[font=\scriptsize,above left=1pt] at (4.55,3.05)
        {$t=\mathrm{cte}$};
    \foreach \x/\time in {-1.95/12{:}02,1.95/11{:}58}{
        \fill[black!75] (\x,3.05) circle (1.8pt);
        \node[font=\scriptsize,below=3pt] at (\x,3.05) {$\time$};
    }
    \node[minkdarkred,font=\small,align=left] at (-3.25,5.35)
        {\contour{white}{sincronizados en $S'$}\\
         \contour{white}{desincronizados para $S$}};
\end{tikzpicture}
\caption{Situación recíproca. Los pequeños relojes de $S'$ están sincronizados sobre la franja roja inclinada; una línea horizontal de simultaneidad de $S$ los intercepta en lecturas distintas.}
\label{fig:relojes_Sp}
\end{figure}